\documentclass[11pt]{article}

\usepackage[margin=1in]{geometry}
\usepackage{amsmath,amssymb}
\usepackage{enumitem}
\usepackage{hyperref}
\usepackage{xcolor}
\usepackage{booktabs}
\usepackage{listings}

\hypersetup{colorlinks=true, linkcolor=blue, urlcolor=blue, citecolor=blue}

\lstset{
  basicstyle=\ttfamily\small,
  breaklines=true,
  columns=fullflexible,
}

\title{TOAE209/TOAH209 -- Design and Analysis of Algorithms\\[4pt]
\large Week 11: Practical Exercises}
\author{Foreign Trade University}
\date{}

\begin{document}
\maketitle

\section*{Instructions}

For each problem below there is a starter file
\texttt{exercises/starter\_code\_problemNN.py} containing function signatures with
\texttt{raise NotImplementedError} bodies, and a small \texttt{\_\_main\_\_} block with
tests. Implement the missing functions so the tests pass. A reference solution is
provided in \texttt{solutions/solution\_problemNN.py} -- try the problem yourself before
looking!

All problems use only the Python 3.10+ standard library. Shared helper code lives in
\texttt{starter\_code.py} at the week root: the \texttt{MaxFlow} class (an Edmonds--Karp
maximum-flow solver with \texttt{add\_edge}, \texttt{max\_flow}, \texttt{flow}, and
\texttt{min\_cut\_reachable}), plus \texttt{generate\_bipartite\_graph},
\texttt{generate\_random\_graph}, \texttt{edges\_to\_adj}, \texttt{matching\_size}, and
\texttt{is\_valid\_matching}. Several later problems reuse earlier solutions (e.g.\
Kuhn's matching from Problem 2), so the starter files add the \texttt{solutions/}
directory to the import path.

The 12 problems are grouped into three themes:
\begin{itemize}
    \item \textbf{Bipartite Matching} (Problems 1--6): bipartite testing,
    maximum matching via both Kuhn's augmenting-path algorithm and the max-flow
    reduction, Hall's deficiency certificate, K\H{o}nig's minimum vertex cover,
    and assignment feasibility.
    \item \textbf{Flow Applications} (Problems 7--10): edge- and vertex-disjoint
    paths, baseball elimination, and project selection via min-cut.
    \item \textbf{Reductions} (Problems 11--12): circulation-with-demands
    feasibility, and verifying the matching/vertex-cover duality (weak in general,
    tight in bipartite graphs).
\end{itemize}

\newpage
\section{Bipartite Matching}

\subsection*{Problem 1 -- Bipartite Check via 2-Coloring}

Implement \texttt{bipartite\_partition(vertices, edges)} that decides whether an
undirected graph is bipartite using a BFS 2-coloring. If it is bipartite, return a tuple
of two sets \texttt{(A, B)} giving the bipartition; if the graph contains an odd cycle,
return \texttt{None}. Run BFS from each unvisited vertex, coloring each neighbour the
opposite color; if an edge ever joins two same-colored vertices, the graph is not
bipartite.

\textbf{Example.} A 4-cycle $0$--$1$--$2$--$3$--$0$ is bipartite (\texttt{\{0,2\}},
\texttt{\{1,3\}}); a triangle is not.

\subsection*{Problem 2 -- Maximum Bipartite Matching via Augmenting Paths (Kuhn's)}

Implement \texttt{kuhn\_matching(left, right, edges)}: maximum bipartite matching by
Kuhn's augmenting-path algorithm. Return a dict mapping each matched \emph{right} vertex
to its \emph{left} partner. For each left vertex, run a DFS that looks for an augmenting
path, bumping already-matched right vertices to alternative partners. Also implement
\texttt{max\_matching\_size(left, right, edges)}.

\subsection*{Problem 3 -- Maximum Bipartite Matching via Max-Flow Reduction}

Implement \texttt{flow\_matching(left, right, edges)} using the \texttt{MaxFlow} helper:
source $\to$ each left vertex (capacity $1$), each edge left $\to$ right (capacity $1$),
each right vertex $\to$ sink (capacity $1$). The saturated left$\to$right edges form a
maximum matching; return them. Also implement \texttt{flow\_matching\_size}. The tests
\textbf{verify the flow-based size equals Kuhn's result} (Problem 2) on random bipartite
graphs -- two very different algorithms computing the same quantity.

\subsection*{Problem 4 -- Hall's Condition Checker}

Implement \texttt{hall\_violator(left, right, edges)}: decide whether the \emph{left}
side has a matching saturating \emph{all} of its vertices. Return \texttt{None} if such a
matching exists; otherwise return a \emph{deficient} subset $S$ of left vertices with
$|N(S)| < |S|$, witnessing the failure of Hall's condition. Recover $S$ from the
alternating-path structure of a maximum matching (Problem 2). Also implement
\texttt{neighbors\_of(subset, adj)}.

\textbf{Hint.} If not all left vertices are matched, the set of left vertices reachable
by alternating paths from the unmatched left vertices is deficient.

\subsection*{Problem 5 -- K\H{o}nig's Theorem: Min Vertex Cover from Max Matching}

Implement \texttt{min\_vertex\_cover(left, right, edges)} via the constructive proof of
K\H{o}nig's theorem: find a maximum matching; from the unmatched left vertices, mark all
vertices reachable by alternating paths; the cover is (unmarked left) $\cup$ (marked
right). Also implement \texttt{is\_vertex\_cover(edges, cover)}. The tests \textbf{verify
$|\text{cover}| = |\text{matching}|$} on random bipartite graphs.

\subsection*{Problem 6 -- Assignment Feasibility (Perfect Matching of One Side)}

Implement \texttt{can\_assign\_everyone(people, tasks, allowed)}: return \texttt{True}
iff every person can be assigned a distinct allowed task (a matching saturating all
people). Also implement \texttt{assignment(people, tasks, allowed)} returning a
\texttt{person -> task} dict if everyone can be assigned, else \texttt{\{\}}. Build on
\texttt{kuhn\_matching} (Problem 2).

\textbf{Example.} Three people each with at least one distinct allowed task can all be
assigned; two people who both only accept a single task cannot.

\newpage
\section{Flow Applications}

\subsection*{Problem 7 -- Maximum Edge-Disjoint s-t Paths}

Implement \texttt{max\_edge\_disjoint\_paths(vertices, directed\_edges, s, t)}: the
maximum number of edge-disjoint directed $s$--$t$ paths. By Menger's theorem this equals
the max flow when every edge has capacity $1$ -- build that network with \texttt{MaxFlow}
and return the flow value.

\textbf{Example.} In a graph with two parallel routes $s\to a\to t$ and $s\to b\to t$,
the answer is $2$; a single bottleneck edge shared by all routes forces the answer down
to $1$.

\subsection*{Problem 8 -- Vertex-Disjoint s-t Paths via Node Splitting}

Implement \texttt{max\_vertex\_disjoint\_paths(vertices, directed\_edges, s, t)}: the
maximum number of internally vertex-disjoint directed $s$--$t$ paths. Use node splitting:
each internal vertex $v$ becomes $v_{\text{in}} \to v_{\text{out}}$ with capacity $1$, and
each original edge $(u,v)$ becomes $u_{\text{out}} \to v_{\text{in}}$ with capacity $1$;
keep $s$ and $t$ uncapped.

\textbf{Hint.} Compare with Problem 7: two paths can be edge-disjoint yet share an
internal vertex (so vertex-disjoint counts can be strictly smaller).

\subsection*{Problem 9 -- Baseball Elimination via Max Flow}

Implement \texttt{is\_eliminated(teams, wins, remaining, games\_between, x)}: decide
whether team \texttt{x} can still possibly finish in first place. Let $W = $
\texttt{wins[x] + remaining[x]}. If any other team already has more wins than $W$, $x$ is
trivially eliminated. Otherwise build the network: source $\to$ game node $\{i,j\}$
(capacity = games left between $i,j$); game node $\to$ each of $i,j$ (capacity $\infty$);
team $i \to$ sink (capacity $W - \texttt{wins[i]}$). Team $x$ survives iff the max flow
\emph{saturates all game edges}.

\textbf{Note.} \texttt{games\_between} is keyed by unordered pairs given as
\texttt{tuple(sorted((i, j), key=str))}.

\subsection*{Problem 10 -- Project Selection / Max-Weight Closure via Min-Cut}

Implement \texttt{select\_projects(profits, prerequisites)}: choose a prerequisite-closed
subset of projects maximizing total profit (profits may be negative). Reduction: source
$\to v$ (capacity \texttt{profit[v]}) for profitable $v$; $v \to$ sink (capacity
$-\texttt{profit[v]}$) for unprofitable $v$; prerequisite edge $p \to q$ with capacity
$\infty$. Then max profit $=$ (sum of positive profits) $-$ (min cut), and the selected
set is the source side (\texttt{MaxFlow.min\_cut\_reachable}). Return
\texttt{(max\_profit, selected\_set)}. Also implement \texttt{brute\_force\_select} (try
all closed subsets) so the tests can \textbf{verify min-cut against brute force} on small
inputs.

\newpage
\section{Reductions}

\subsection*{Problem 11 -- Circulation with Demands: Feasibility Check}

Implement \texttt{circulation\_feasible(vertices, demands, edges)}: decide whether a
feasible circulation with demands and lower bounds exists. Each vertex $v$ has a demand
$d_v$ ($>0$ absorb, $<0$ supply); each edge \texttt{(u, v, lower, upper)} must carry flow
in $[\texttt{lower}, \texttt{upper}]$. Reduce to max flow: (1) eliminate lower bounds by
setting capacity to \texttt{upper - lower} and adjusting $d_u \mathrel{+}= \texttt{lower}$,
$d_v \mathrel{-}= \texttt{lower}$; (2) add a super-source/super-sink for the adjusted
supplies/demands; (3) feasible iff adjusted supply equals adjusted demand \emph{and} the
max flow saturates all super-source edges.

\textbf{Example.} A source $a$ ($-3$) feeding a sink $b$ ($+3$) over an edge of capacity
$\ge 3$ is feasible; the same with capacity $2$, or with a lower bound that forces more
flow than $b$ can absorb, is infeasible.

\subsection*{Problem 12 -- Matching $\le$ Vertex Cover, with Equality in Bipartite Graphs}

Implement \texttt{brute\_force\_max\_matching(vertices, edges)} and
\texttt{brute\_force\_min\_vertex\_cover(vertices, edges)} (both for \emph{any} graph, by
brute force), and \texttt{verify\_matching\_cover\_relationship(num\_trials, seed)} which,
on random graphs, checks: (1) maximum matching $\le$ minimum vertex cover always (weak
duality); and (2) in \emph{bipartite} graphs maximum matching $=$ minimum vertex cover
(K\H{o}nig), with the constructive cover from Problem 5 attaining the bound.

\textbf{Example.} A triangle (non-bipartite) has maximum matching $1$ but minimum vertex
cover $2$ -- a strict inequality, confirming K\H{o}nig's theorem requires bipartiteness.

\newpage
\section{LeetCode Practice}

After completing the problems above, practise this week's flow and matching techniques
(bipartite checking, bipartite matching, and graph/DAG modeling) on the following
\href{https://leetcode.com}{LeetCode} problems. Difficulty is LeetCode's own label.

\begin{itemize}
  \item \href{https://leetcode.com/problems/is-graph-bipartite/}{Is Graph Bipartite?} -- \emph{Medium} -- Bipartite check.
  \item \href{https://leetcode.com/problems/possible-bipartition/}{Possible Bipartition} -- \emph{Medium} -- Bipartite check.
  \item \href{https://leetcode.com/problems/maximum-number-of-accepted-invitations/}{Maximum Number of Accepted Invitations} -- \emph{Medium} -- Bipartite matching.
  \item \href{https://leetcode.com/problems/optimal-account-balancing/}{Optimal Account Balancing} -- \emph{Hard} -- Flow / search.
  \item \href{https://leetcode.com/problems/minimum-number-of-vertices-to-reach-all-nodes/}{Minimum Number of Vertices to Reach All Nodes} -- \emph{Medium} -- Graph application.
  \item \href{https://leetcode.com/problems/course-schedule-ii/}{Course Schedule II} -- \emph{Medium} -- DAG application.
\end{itemize}

\end{document}
